% III. Експериментална проверка и резултати.
% Всяко число тук е измерено на посочената машина. Суровият изход на трите
% изпълнения е запазен в docs/measurements/. Таблиците са от второто изпълнение.
\section{Експериментална проверка и резултати}
\label{sec:method}
\label{sec:results}

\subsection{Коректност}

Синтаксисът се проверява със съставен за целта набор от 101 случая: 41 за Turtle (26
положителни и 15 отрицателни), 16 за N-Triples (7 и 9) и 44 за SPARQL (26 и 18). Всички
преминават. Отрицателният случай се проверява два пъти: веднъж, че изобщо е отхвърлен, и втори
път, че отказът носи ред, колона и текст, защото анализатор, който казва само „синтактична
грешка“, не върши работа. Осемнадесетте отрицателни случая за SPARQL са конструкциите извън
подмножеството, тоест проверени граници, а не дефекти. Наборът е съставен в хранилището, а не
изтеглен, защото тестовите набори на W3C~\cite{w3c_sparql11_testsuite} са изтегляне, а проектът
трябва да се изгражда без мрежов достъп; затова тези числа не бива да се четат като резултат от
нормативния набор.

Семантиката се проверява върху малък граф от 40 триплета, чиито отговори са преброени на ръка
от текста на файла, а не отчетени от реализацията. Пълният набор съдържа 125 случая в 10 групи
в CTest и всичките преминават. Отделно се проверяват и двете твърдения на
Раздел~\ref{sec:design}: че за всяка от осемте комбинации избраният индекс има префикс с
дължина, равна на броя свързани позиции, и че планираният ред, наивният ред, сливането и
вложените цикли дават едно множество решения.

\subsection{Условия за валидност на измерването}

Наборът се генерира от хранилището и описва библиографски граф: статии със заглавие, година,
конференция, тема, брой цитирания и автори; автори с име, месторабота и показател, като само
част от тях имат домашна страница, за да има какво да показва незадължителното съединение.
Схемната част добавя подкласове, подсвойство, област и обхват, а работното натоварване е осемте
заявки от приложение~В.

Измерванията са проведени на AMD Ryzen 5 3600 с 6 ядра и 15{,}9 GiB памет, под Windows 11
Pro N, с g++ 15.2.0 (MinGW-w64) в режим \code{CMAKE\_BUILD\_TYPE=Release}, CMake 4.3.2 и Ninja
1.13.2. Всяка конфигурация се изпълнява 15 пъти, първите две се отхвърлят заради загряването на
разпределителя на памет, и се докладва медианата. Броят на решенията се отпечатва до всяко
време: ред, в който броевете не съвпадат, не е сравнение на скорост. Машината е работна, а не
изолиран стенд, затова наборът е повторен три пъти; суровият изход е в
\code{docs/measurements/}, а таблиците са от второто изпълнение.

\subsection{Резултати: зареждане и изпълнение на заявки}

\begin{table}[H]
\caption{Зареждане, изграждане на индексите и обем}
\label{tab:loading}
\centering
\fontsize{9}{11}\selectfont
\setlength{\tabcolsep}{4pt}
\begin{tabular}{@{}rrrrrrr@{}}
\toprule
\textbf{Статии} & \textbf{Триплети} & \textbf{Вход, B} & \textbf{Анализ, ms} & \textbf{Триплети/s} & \textbf{Индекси, ms} & \textbf{Обем, KiB} \\
\midrule
1\,000  &  10\,462 &     301\,893 &  12{,}2 &   858\,239 &   2{,}0 &  1\,827 \\
5\,000  &  52\,395 &  1\,528\,499 &  59{,}4 &   882\,778 &  11{,}1 &  7\,024 \\
20\,000 & 209\,202 &  6\,185\,046 & 261{,}4 &   800\,498 &  57{,}4 & 28\,032 \\
50\,000 & 522\,910 & 15\,546\,537 & 504{,}6 & 1\,036\,328 & 114{,}5 & 72\,123 \\
\bottomrule
\end{tabular}
\end{table}

Броят на различните термини в речника е 2\,135, 7\,614, 27\,565 и 67\,501. Времето за анализ
расте линейно \tabref{tab:loading}, а пропускателната способност остава от порядъка на
800\,000 до 1\,000\,000 триплета за секунда, без спад при петдесеткратно нарастване на входа.
Изграждането на трите индекса е около една пета от времето за анализ, тоест сортиранията не са
водещият разход, а обемът е около 141 байта на триплет, което включва трите копия по 24 байта
и речника. Разсейването е най-голямо именно тук: при 20\,000 статии способността е 474\,237,
800\,498 и 905\,522 триплета за секунда, тоест почти два пъти разлика, и величината следва да
се чете като порядък.

\begin{table}[H]
\caption{Медианно време при включен и при изключен планировчик, 209\,202 триплета}
\label{tab:queries}
\centering
\fontsize{10}{12}\selectfont
\begin{tabular}{@{}L{5.4cm}rrrr@{}}
\toprule
\textbf{Заявка} & \textbf{Планирано, ms} & \textbf{Наивно, ms} & \textbf{Отношение} & \textbf{Решения} \\
\midrule
B1 един шаблон                & 1{,}375 & 1{,}400 & 1{,}02 & 20\,000 \\
B2 звезда от два шаблона      & 6{,}822 & 6{,}831 & 1{,}00 & 20\,000 \\
B3 верига от четири шаблона   & 0{,}028 & 12{,}705 & 447{,}4 & 47 \\
B4 верига от пет шаблона      & 0{,}028 & 12{,}833 & 465{,}0 & 0 \\
B5 филтър по число            & 6{,}172 & 6{,}122 & 0{,}99 & 3\,188 \\
B6 незадължително съединение  & 1{,}352 & 1{,}348 & 1{,}00 & 4\,000 \\
B7 групиране с брояч          & 2{,}613 & 2{,}918 & 1{,}12 & 200 \\
B8 подреждане на цялото       & 13{,}623 & 13{,}894 & 1{,}02 & 20\,000 \\
\bottomrule
\end{tabular}
\end{table}

Колоната с броя решения \tabref{tab:queries} е условие за валидност: двата режима връщат едно и
също множество при всяка заявка. Нулата при B4 е действителният отговор, а не пропуснато измерване, тъй като
авторите от избраната организация нямат статия в избраната конференция; редът остава, защото
времето до празния отговор се различава 465 пъти.

Стратегията на съединяване беше сравнена върху заявката с два шаблона със свързано сказуемо,
споделящи допълнението си, тоест единствената форма, при която сливането е допустимо. Двете
връщат по 99\,996 решения; сливането отнема 10{,}313, 8{,}995 и 9{,}149 ms, а вложените цикли
по индекс 10{,}112, 10{,}049 и 10{,}042 ms. При мощност на водещия шаблон 4, 5 и 11 заявката
отнема 0{,}022, 0{,}022 и 0{,}028 ms и връща 29, 30 и 47 решения, тоест времето расте с
мощността, а не с размера на набора. Материализацията по RDFS извежда 68\,398 нови триплета за
две обиколки до неподвижна точка, при което наборът става 277\,600, а обемът расте от 28\,032
на 38\,985 KiB за 168{,}6, 312{,}3 и 157{,}5 ms.

\subsection{Анализ}

\textbf{Планирането има смисъл там, където има какво да се планира.} При един или два шаблона
двата реда се различават в границите на шума, а при вериги от четири и пет шаблона отношението
е между 447 и 465 пъти. Причината е пряка: при наивен ред първият шаблон е
\code{?p ex:title ?t} с мощност 20\,000, а при планиран ред първи е
\code{?a ex:affiliation ex:org3} с мощност 11, тоест междинният резултат тръгва с три порядъка
по-малък. Разходът за планиране е под една стотна от милисекундата, а точната мощност се оказа
и по-проста, и по-надеждна от оценка: структурата на индекса вече е статистиката.

\textbf{Стратегията на съединяване не се оказа водеща.} Разликата между сливане и вложени цикли
е в границите на разсейването. Това противоречи на очакването при проектирането и е записано
такова, каквото е; числата допускат обяснението, че при план с лява дълбочина сливането е
приложимо само за първата двойка шаблони.

\textbf{Какво не е измерено.} Сравнение с Apache Jena~\cite{jena_docs} и
RDFLib~\cite{rdflib_docs} беше предвидено, но не е проведено: и двете изискват изтегляне и среда
за изпълнение, каквито на машината не са налични. Число, което не е измерено, няма място в
отчет; работното натоварване и генераторът обаче са записани така, че сравнението да може да
бъде проведено по-късно. Не са измерени върховата резидентна памет на процеса, поведението
при набор, който не се събира в паметта, и цената на извеждане по време на заявка спрямо
материализацията, въпреки че вторият режим вече е реализиран. Времената преди двете промени в
изпълнението също не се докладват: междинните състояния не са запазени. Изводите важат за една
машина.
